% ======================================================================
% INTRODUCTION (4–5 paragraphs)
% P1 field need -> P2 limits of existing methods -> P3 why cilia ->
% P4 contributions -> P5 significance.
% Citations are placeholders (\refneeded) — do NOT invent references.
% ======================================================================
\section*{Introduction}

Embodied intelligence, whether in robotic manipulators or on-body wearable
interfaces, depends on tactile sensing that is at once large-area, mechanically
compliant, multi-axis, dynamic and sensitive to weak forces
\cite{dahiya2010tactile,chortos2016pursuing}. Dexterous manipulation and safe
human contact demand resolving not only firm grasping loads but also light
touch, shear, incipient slip, faint airflow and physiological micro-vibrations
such as the pulse \cite{sundaram2019learning}. Yet most tactile platforms are
optimized for a narrow slice of this space, and no single skin resolves subtle,
spatially distributed and rapidly varying stimuli across an entire hand.

No single transduction principle spans this requirement space. Capacitive and
piezoresistive skins scale to large areas but contend with dense interconnects,
cross-talk, baseline drift and poor weak-force resolution
\cite{sundaram2019learning,boutry2018hierarchically}. Piezoelectric and
triboelectric devices capture dynamic transients but fade under sustained static
load \refneeded{piezoelectric tactile}. Vision-based tactile sensors deliver
fine spatial resolution but demand an internal camera, illumination and
per-contact computation, constraining thickness, curvature and integration
density \cite{yuan2017gelsight}. Conventional magnetic skins offer contactless,
drift-tolerant, multi-axis readout, yet their planar magnet-film geometries
often limit weak-stimulus sensitivity, directional decoupling or scalable
coverage \cite{hellebrekers2019soft,yan2021soft}.

Biology resolves weak mechanical signals with slender, high-aspect-ratio
appendages---mammalian whiskers, insect sensilla and hair cells---that convert
faint airflow, contact and vibration into amplified deflection
\cite{park2023bioinspired}. Magnetic cilia borrow this geometry but have so far
been demonstrated only as single elements or small arrays
\cite{alfadhel2015magnetic,man2023tactile}. We translate the principle into an
ordered, full-hand array of magnetized elastomer cilia. Their high aspect ratio
yields compliance to weak shear, flow and vibration; because each cilium is
transduced magnetically by an underlying tri-axis sensor, the array needs no
individual wiring and forms a passive mechanical amplification layer compatible
with sparse, distributed magnetic sampling.

Here we report a magnetic cilia tactile skin and test it against five questions.
Why cilia, not a flat magnetic film? Can a local $2\times2$ tri-axis unit invert
contact position, force and direction at unseen positions? Does that capability
scale to a 44-node full-hand array? Do the resulting whole-hand maps track true
contact rather than posture, geomagnetic or cable-motion artifacts? And can the
platform capture weak dynamic disturbances---airflow and pulse---that rigid force
sensors miss? Our contributions are (i) a geometry-optimized magnetic cilia
transduction layer with a quantified enhancement over matched flat controls,
(ii) local $2\times2$ inverse decoding validated at unseen positions,
(iii) a 44-node full-hand array with explicit motion-artifact suppression and
trustworthy spatial mapping, and (iv) shared sensing of weak external airflow
and internal pulse micro-vibrations on one platform.

Together, these results establish magnetic cilia as a scalable mechanical
transduction layer for embodied perception---one compliant skin that resolves a
robot's grasp, the airflow around it and the pulse beneath it---pointing toward
integrated robotic touch, wearable interfaces and unobtrusive monitoring of
subtle environmental and physiological signals \refneeded{embodied sensing
outlook}.
